% !TeX program = xelatex
% !TeX encoding = UTF-8
% =========================================================
% hustreport.cls 源代码完整分析教程
% 版本: v3.0
% 语言: 中文
% =========================================================
\documentclass[linux]{hustreport}

% 文档基本信息
\title{hustreport.cls 源代码分析}
\author{HUST Report 开发团队}

\begin{document}

% =========================================================
% 封面页
% =========================================================
\begin{titlepage}
    \centering
    \vspace*{3cm}
    {\Huge\bfseries\heiti hustreport.cls\par}
    \vspace{0.5cm}
    {\Large\songti 黑龙江科技大学报告模板类\par}
    \vspace{1cm}
    {\large 源代码完整分析教程\par}
    \vspace{2cm}
    {\large 版本 v3.0\par}
    \vspace{0.5cm}
    {\large \today\par}
    \vfill
    {\small 本文档对 hustreport.cls 类文件进行逐模块的详细分析。\par}
\end{titlepage}

\tableofcontents
\newpage

% =========================================================
% 第一章: 引言与概述
% =========================================================
\section{引言与概述}

\subsection{模板用途}

\texttt{hustreport.cls} 是为黑龙江科技大学本科生设计的 \LaTeX{} 文档类，主要用于：

\begin{itemize}
    \item 社会调查报告
    \item 实习报告
    \item 毕业论文（使用 \texttt{thesis} 选项）
\end{itemize}

\subsection{架构概览}

类文件包含 \textbf{20 个功能模块}，共约 1445 行代码：

\begin{enumerate}
    \item \textbf{模块 0-1}：版本声明与类选项
    \item \textbf{模块 2}：宏包加载（按功能分类）
    \item \textbf{模块 3}：用户信息接口
    \item \textbf{模块 4}：智能字体配置
    \item \textbf{模块 5-6}：页面布局与章节标题
    \item \textbf{模块 7-8}：图表格式与 PDF 属性
    \item \textbf{模块 9-10}：特殊环境与封面生成
    \item \textbf{模块 11-20}：高级功能（代码块、批注系统等）
\end{enumerate}

\subsection{编译要求}

本模板\textbf{必须}使用 XeLaTeX 编译：

\begin{verbatim}
xelatex your_document.tex
bibtex your_document      % 如有参考文献
xelatex your_document.tex
xelatex your_document.tex
\end{verbatim}

% =========================================================
% 第二章: 模块 0-1 - 版本声明与类选项
% =========================================================
\newpage
\section{模块 0-1：版本声明与类选项}

\subsection{版本声明（第 1-14 行）}

每个 \LaTeX{} 类文件必须声明格式要求并提供标识信息：

\begin{verbatim}
% 第 12-13 行：格式与类标识
\NeedsTeXFormat{LaTeX2e}
\ProvidesClass{hustreport}[2026/01/05 v3.0 HUST Report Class]
\end{verbatim}

\textbf{要点说明：}
\begin{itemize}
    \item \verb|\NeedsTeXFormat{LaTeX2e}| 确保与现代 \LaTeX{} 兼容
    \item \verb|\ProvidesClass| 提供版本信息，编译日志中可见
\end{itemize}

\subsection{布尔开关（第 19-24 行）}

类使用布尔开关控制条件编译：

\begin{verbatim}
% 平台检测开关
\newif\ifhust@oswin    \hust@oswinfalse
\newif\ifhust@osmac    \hust@osmacfalse
\newif\ifhust@oslinux  \hust@oslinuxfalse
\newif\ifhust@optfonts \hust@optfontsfalse
\newif\ifhust@thesis   \hust@thesisfalse
\end{verbatim}

\textbf{实现模式：}
\begin{itemize}
    \item \verb|\newif| 创建布尔开关，默认为 \texttt{false}
    \item 命名约定：\verb|\ifhust@<名称>| 使用内部前缀 \verb|hust@|
    \item 所有开关初始化为 \texttt{false}，采用"显式启用"模式
\end{itemize}

\subsection{选项声明（第 27-31 行）}

使用 \verb|\DeclareOption| 声明选项：

\begin{verbatim}
% 平台选项
\DeclareOption{win}{\hust@oswintrue}
\DeclareOption{mac}{\hust@osmactrue}
\DeclareOption{linux}{\hust@oslinuxtrue}
\DeclareOption{fonts}{\hust@optfontstrue}
\DeclareOption{thesis}{\hust@thesistrue}
\end{verbatim}

\subsection{选项转发与处理（第 33-35 行）}

未识别的选项转发给基类：

\begin{verbatim}
\DeclareOption*{\PassOptionsToClass{\CurrentOption}{ctexrep}}
\ProcessOptions\relax
\end{verbatim}

\subsection{基类加载（第 37-43 行）}

模板基于 \texttt{ctexrep} 构建以支持中文：

\begin{verbatim}
\LoadClass[a4paper,UTF8,zihao=-4,linespread=1.5]{ctexrep}
\end{verbatim}

\textbf{参数说明：}
\begin{optionlist}
    \option{a4paper}{标准 A4 纸张尺寸（210mm × 297mm）}
    \option{UTF8}{UTF-8 编码支持中文}
    \option{zihao=-4}{默认字号：小四（相当于 12pt）}
    \option{linespread=1.5}{1.5 倍行距提高可读性}
\end{optionlist}

% =========================================================
% 第三章: 模块 2 - 宏包加载
% =========================================================
\newpage
\section{模块 2：宏包加载}

\subsection{组织理念}

宏包按功能类别组织以便维护：

\subsection{页面与布局宏包（第 51-55 行）}

\begin{verbatim}
\RequirePackage{geometry}   % 页面边距
\RequirePackage{fancyhdr}   % 页眉页脚
\RequirePackage{float}      % 浮动体控制
\RequirePackage[sort&compress,numbers,square]{natbib}
\end{verbatim}

\subsection{标题与目录宏包（第 57-59 行）}

\begin{verbatim}
\RequirePackage{titletoc}   % 目录格式
\RequirePackage{titlesec}   % 标题格式
\end{verbatim}

\subsection{数学与物理宏包（第 61-66 行）}

\begin{verbatim}
\RequirePackage{amsmath,amssymb,amsfonts,amsthm}
\RequirePackage{physics}    % 物理符号简写
\RequirePackage{siunitx}    % 国际单位制
\RequirePackage{newtxmath}  % 数学字体优化
\RequirePackage{calc}       % 长度计算
\end{verbatim}

\subsection{图形与颜色宏包（第 68-72 行）}

\begin{verbatim}
\RequirePackage{graphicx}   % 图片插入
\RequirePackage{subcaption} % 子图支持
\RequirePackage{tikz}       % 绘图与定位
\RequirePackage[dvipsnames,svgnames,x11names]{xcolor}
\end{verbatim}

\subsection{表格与列表宏包（第 74-84 行）}

\begin{verbatim}
\RequirePackage{booktabs}   % 专业表格
\RequirePackage{threeparttable}
\RequirePackage{caption}
\RequirePackage{tabularx}   % 自适应宽度表格
\RequirePackage{xltabular}  % 跨页长表格
\RequirePackage{multirow}   % 多行合并
\RequirePackage{makecell}   % 单元格换行
\RequirePackage{array}      % 列格式扩展
\RequirePackage{enumitem}   % 列表定制
\RequirePackage{diagbox}    % 斜线表头
\end{verbatim}

% =========================================================
% 第四章: 模块 3 - 信息接口
% =========================================================
\newpage
\section{模块 3：信息接口}

\subsection{内部存储变量（第 97-102 行）}

类定义使用 \verb|@| 前缀的内部变量：

\begin{verbatim}
\makeatletter
\newcommand{\@major}{}      % 专业
\newcommand{\@dept}{}       % 院系
\newcommand{\@advisor}{}    % 指导教师
\newcommand{\@team}{}       % 队员
\newcommand{\@studentID}{}  % 学号
\newcommand{\@reportID}{}   % 报告编号
\end{verbatim}

\subsection{对外接口命令（第 105-110 行）}

用户使用的命令，用于设置内部变量：

\begin{verbatim}
\newcommand{\major}[1]{\def\@major{#1}}
\newcommand{\dept}[1]{\def\@dept{#1}}
\newcommand{\advisor}[1]{\def\@advisor{#1}}
\newcommand{\team}[1]{\def\@team{#1}}
\newcommand{\studentID}[1]{\def\@studentID{#1}}
\newcommand{\reportID}[1]{\def\@reportID{#1}}
\end{verbatim}

\subsection{任务书变量（第 113-116 行）}

任务书内容的预定义变量：

\begin{verbatim}
\newcommand{\taskTargetText}{内容待定}
\newcommand{\taskContentText}{内容待定}
\newcommand{\taskExpectedText}{内容待定}
\newcommand{\taskScheduleText}{进度待定}
\end{verbatim}

\subsection{智能下划线组件（第 119-123 行）}

表单填空处的自适应下划线：

\begin{verbatim}
\newcommand{\ulineinfo}[2][280pt]{%
    \begin{minipage}[t]{#1}%
        \centering\zihao{3}\songti\uline{#2\hfill}%
    \end{minipage}%
}
\end{verbatim}

% =========================================================
% 第五章: 模块 4 - 字体配置
% =========================================================
\newpage
\section{模块 4：智能字体配置}

\subsection{设计理念}

字体系统实现\textbf{回退链}，包含三个优先级：

\begin{enumerate}
    \item \textbf{最高优先级}：本地 \texttt{fonts/} 目录
    \item \textbf{中等优先级}：平台特定系统字体
    \item \textbf{保底策略}：Fandol 开源字体
\end{enumerate}

\subsection{字体重置（第 129-131 行）}

清除现有字体定义以防冲突：

\begin{verbatim}
\let\songti\relax
\let\heiti\relax
\let\kaiti\relax
\end{verbatim}

\subsection{英文字体配置（第 147-167 行）}

\begin{verbatim}
\IfFileExists{fonts/times.ttf}{
    % 找到本地 Times 字体
    \setmainfont{times.ttf}[
        Path = fonts/,
        BoldFont = timesbd.ttf,
        ItalicFont = timesi.ttf,
        BoldItalicFont = timesbi.ttf
    ]
}{
    % 系统回退
    \ifhust@oswin
        \setmainfont{Times New Roman}
    \else
        \ifhust@osmac
            \setmainfont{Times New Roman}
        \else
            \setmainfont{TeX Gyre Termes}
        \fi
    \fi
}
\end{verbatim}

\subsection{中文字体配置（第 171-254 行）}

中文字体设置使用条件检查每个字体文件：

\begin{verbatim}
\IfFileExists{fonts/SimSun.ttf}{
    % 本地中文字体可用
    \setCJKmainfont{SimSun.ttf}[
        Path = fonts/,
        BoldFont = \hust@cjk@bold,
        ItalicFont = \hust@cjk@italic,
        AutoFakeBold = \hust@cjk@fakebold
    ]
    % ... 其他字体族定义
}{
    % 平台特定系统字体
    \ifhust@oswin
        % Windows: SimSun, SimHei, KaiTi
    \else\ifhust@osmac
        % macOS: Songti SC, Heiti SC, Kaiti SC
    \else
        % Linux: Fandol 字体
    \fi\fi
}
\end{verbatim}

% =========================================================
% 第六章: 模块 5-6 - 页面布局与章节标题
% =========================================================
\newpage
\section{模块 5-6：页面布局与章节标题}

\subsection{条件页面布局（第 260-301 行）}

布局根据 \texttt{thesis} 选项切换：

\begin{verbatim}
\ifhust@thesis
    % 毕业论文模式：带校徽的特殊页眉
    \geometry{
        a4paper,
        top=3.7cm, bottom=2.5cm,
        left=3.9cm, right=2.5cm,
        headheight=3cm, headsep=0.5cm
    }
    
    % 带校徽的自定义页眉
    \pagestyle{fancy}
    \fancyhf{}
    \renewcommand{\headrulewidth}{2.5pt}
    \fancyhead[L]{\includegraphics[height=2.6cm]{logo.png}}
    \fancyhead[C]{...校名和标题...}
    \fancyfoot[C]{\thepage}
\else
    % 标准报告模式
    \geometry{left=3.017cm, right=3.017cm, 
              top=2.54cm, bottom=2.54cm}
    \pagestyle{plain}
\fi
\end{verbatim}

\subsection{计数器配置（第 305-313 行）}

章节编号使用阿拉伯数字并与章绑定：

\begin{verbatim}
\renewcommand{\thesection}{\arabic{section}}
\renewcommand{\thesubsection}{\thesection.\arabic{subsection}}
\renewcommand{\thesubsubsection}{\thesubsection.\arabic{subsubsection}}

% 计数器绑定到章节（每章重置）
\counterwithin{equation}{section}
\counterwithin{figure}{section}
\counterwithin{table}{section}
\end{verbatim}

\subsection{章节标题样式（第 325-352 行）}

使用 \texttt{ctexset} 进行中文感知格式化：

\begin{verbatim}
\ctexset{
    section = {
        number = \chinese{section},
        name = {第,章},
        format = \centering\zihao{3}\heiti\bfseries,
        aftername = \quad,
        beforeskip = 15pt, afterskip = 20pt
    },
    subsection = {
        format = \zihao{4}\heiti\bfseries,
        aftername = \hspace{1em}
    },
    subsubsection = {
        format = \zihao{-4}\songti,
        aftername = \hspace{1em}
    }
}
\end{verbatim}

\subsection{目录样式（第 358-376 行）}

自定义目录格式：

\begin{verbatim}
\titlecontents{section}[0pt]
{\addvspace{6pt}\bfseries\heiti\zihao{4}}
{\thecontentslabel\quad}
{}
{\titlerule*[0.5pc]{.}\contentspage}
\end{verbatim}

% =========================================================
% 第七章: 模块 7-8 - 图表与 PDF
% =========================================================
\newpage
\section{模块 7-8：图表格式与 PDF 属性}

\subsection{浮动体参数优化（第 386-392 行）}

解决"图片乱跑"问题：

\begin{verbatim}
\renewcommand{\topfraction}{0.85}    % 顶部浮动体：85%
\renewcommand{\bottomfraction}{0.70} % 底部浮动体：70%
\renewcommand{\textfraction}{0.10}   % 最小文本：10%
\renewcommand{\floatpagefraction}{0.50}
\end{verbatim}

\subsection{标题样式（第 395-410 行）}

专业的标题格式：

\begin{verbatim}
\DeclareCaptionFont{heiti}{\heiti}
\DeclareCaptionFont{fivehao}{\zihao{5}}
\captionsetup{
    format=hang,
    labelsep=quad,
    font={fivehao,heiti},
    skip=6pt
}
\end{verbatim}

\subsection{编号格式（第 415-423 行）}

相对于章节的编号：

\begin{verbatim}
\renewcommand{\theequation}{\thesection-\arabic{equation}}
\renewcommand{\thefigure}{\thesection-\arabic{figure}}
\renewcommand{\thetable}{\thesection-\arabic{table}}

% 引用自动添加前缀
\labelformat{figure}{图~#1}
\labelformat{table}{表~#1}
\labelformat{equation}{式~(#1)}
\end{verbatim}

\subsection{PDF 属性（第 442-448 行）}

Hyperref 配置：

\begin{verbatim}
\hypersetup{
    hidelinks,
    colorlinks=false,
    bookmarksnumbered=true,
    bookmarksopen=true,
    pdfstartview=FitH
}
\end{verbatim}

% =========================================================
% 第八章: 模块 9-10 - 特殊环境与封面
% =========================================================
\newpage
\section{模块 9-10：特殊环境与封面}

\subsection{中文摘要环境（第 455-468 行）}

\begin{verbatim}
\newenvironment{cnabstract}[1]{
    \def\@keyword{#1}
    \clearpage
    \phantomsection
    \addcontentsline{toc}{section}{摘要}
    \begin{center} \heiti\zihao{3} \@title \end{center}
    \vspace{1ex}
    \begin{center} \heiti\zihao{4} 摘　要 \end{center}
    \par
}{
    \par\vspace{2ex}
    \noindent{\heiti\zihao{-4} 关键词：}\@keyword
    \newpage
}
\end{verbatim}

\subsection{参考文献自动化（第 471-498 行）}

\begin{verbatim}
\newcommand{\MyBibliography}[1]{
    \clearpage
    \phantomsection
    \addcontentsline{toc}{section}{参考文献}
    \renewcommand{\section}[2]{}% 抑制自动标题
    \bibliographystyle{plain}
    \bibliography{#1}
}
\end{verbatim}

\subsection{附录环境（第 501-527 行）}

\begin{verbatim}
\newenvironment{appendixblock}{
    \clearpage
    \setcounter{section}{0}
    \renewcommand{\thesection}{\Alph{section}}
    \ctexset{
        section = {
            name = {附录,},
            format = \centering\zihao{4}\heiti\bfseries
        }
    }
}{}
\end{verbatim}

\subsection{TikZ 封面生成（第 641-680 行）}

封面使用 TikZ 进行绝对定位：

\begin{verbatim}
\newcommand{\makecover}{
    \begin{titlepage}
        \begin{tikzpicture}[remember picture, overlay]
            % 校徽定位
            \node[anchor=north east, xshift=-2cm, yshift=-2cm] 
                at (current page.north east) {
                \includegraphics[width=0.25\textwidth]{logo.png}
            };
        \end{tikzpicture}
        % 标题和信息表格...
    \end{titlepage}
}
\end{verbatim}

% =========================================================
% 第九章: 模块 11-12 - 高级布局与代码
% =========================================================
\newpage
\section{模块 11-12：高级布局与代码环境}

\subsection{图文混排布局（第 693-737 行）}

\texttt{textfigure} 环境使用 \texttt{xparse}：

\begin{verbatim}
\NewDocumentEnvironment{textfigure}{O{right} m m O{0.4}}{%
    \def\tf@pos{#1}%    % 位置：left 或 right
    \def\tf@path{#2}%   % 图片路径
    \def\tf@cap{#3}%    % 标题
    \setlength{\tf@imagew}{#4\textwidth}%
    % 计算文字宽度
    \setlength{\tf@textw}{\textwidth}%
    \addtolength{\tf@textw}{-\tf@imagew}%
    \addtolength{\tf@textw}{-1.5em}%
    % 开始文字捕获...
}{%
    % 并排渲染图片 + 文字
}
\end{verbatim}

\subsection{代码环境设计（第 772-811 行）}

语法高亮的颜色定义：

\begin{verbatim}
\definecolor{hust@codeback}{RGB}{250,250,250}
\definecolor{hust@keyword}{RGB}{0,92,197}
\definecolor{hust@comment}{RGB}{80,161,79}
\definecolor{hust@string}{RGB}{227,17,108}

\lstdefinestyle{huststyle}{
    basicstyle=\small\ttfamily,
    keywordstyle=\bfseries\color{hust@keyword},
    commentstyle=\itshape\color{hust@comment},
    stringstyle=\color{hust@string},
    numbers=left,
    breaklines=true,
    keepspaces=true
}
\end{verbatim}

\subsection{tcolorbox 集成（第 814-855 行）}

Mac 风格三色点代码块：

\begin{verbatim}
\tcbset{
    hustcodestyle/.style={
        enhanced, breakable,
        colback=hust@codeback,
        title={\textbf{代码 \thehustcode：}#1},
        % Mac 风格装饰
        overlay unbroken and first={
            \draw[fill=red!70!white] 
                ([xshift=6mm,yshift=-1.5mm]frame.north west) 
                circle (2.5pt);
            \draw[fill=yellow!70!white] 
                ([xshift=11mm,yshift=-1.5mm]frame.north west) 
                circle (2.5pt);
            \draw[fill=green!70!white] 
                ([xshift=16mm,yshift=-1.5mm]frame.north west) 
                circle (2.5pt);
        }
    }
}
\end{verbatim}

% =========================================================
% 第十章: 模块 13-16 - 摘要、算法、提示框
% =========================================================
\newpage
\section{模块 13-16：摘要、算法与提示框环境}

\subsection{英文摘要（第 872-887 行）}

\begin{verbatim}
\newenvironment{enabstract}[1]{
    \def\@enkeyword{#1}
    \clearpage
    \addcontentsline{toc}{section}{Abstract}
    \begin{center}\bfseries\zihao{4} Abstract \end{center}
    \linespread{1.5}\selectfont
    \zihao{-4}
}{
    \par\vspace{2ex}
    \noindent\textbf{Keywords: } \@enkeyword
    \newpage
}
\end{verbatim}

\subsection{算法汉化（第 892-904 行）}

\begin{verbatim}
\RequirePackage{algorithm}
\RequirePackage{algorithmicx}
\RequirePackage{algpseudocode}

\floatname{algorithm}{算法}
\renewcommand{\algorithmicrequire}{\textbf{输入：}}
\renewcommand{\algorithmicensure}{\textbf{输出：}}
\algrenewcommand\algorithmicindent{2.0em}
\end{verbatim}

\subsection{提示框（第 928-957 行）}

使用 tcolorbox 的三种主题框：

\begin{verbatim}
% 注意框（蓝色主题）
\newtcolorbox{notice}[1]{
    commonbox,
    colback=Azure1, colframe=DodgerBlue3,
    title={【注意】#1}
}

% 定义框（绿色主题）
\newtcolorbox{definition}[1]{
    commonbox,
    colback=Honeydew1, colframe=ForestGreen,
    title={【定义】#1}
}

% 结论框（红色主题）
\newtcolorbox{conclusion}[1]{
    commonbox,
    colback=LavenderBlush1, colframe=Crimson,
    title={【结论】#1}
}
\end{verbatim}

\subsection{致谢环境（第 962-970 行）}

\begin{verbatim}
\newenvironment{acknowledgement}{
    \clearpage
    \addcontentsline{toc}{section}{致谢}
    \section*{致谢}
    \zihao{-4}\songti\linespread{1.5}\selectfont
}{
    \par
}
\end{verbatim}

% =========================================================
% 第十一章: 模块 17 - 列表定制
% =========================================================
\newpage
\section{模块 17：专业列表定制}

\subsection{自定义 TikZ 列表标记（第 981-1005 行）}

\begin{verbatim}
% 一级：实心方块
\newcommand{\hust@item@square}{%
    \tikz[baseline=-0.6ex] 
    \node[fill=black, shape=rectangle, 
          inner sep=0pt, minimum size=0.5em] {};%
}

% 二级：灰色圆点
\newcommand{\hust@item@circle}{%
    \tikz[baseline=-0.6ex] 
    \node[fill=darkgray, shape=circle, 
          inner sep=0pt, minimum size=0.35em] {};%
}

% 三级：短横线
\newcommand{\hust@item@dash}{%
    \tikz[baseline=-0.6ex] 
    \draw[line width=0.8pt] (0,0) -- (0.5em,0);%
}
\end{verbatim}

\subsection{无序列表配置（第 1018-1031 行）}

\begin{verbatim}
\setlist[itemize,1]{
    label=\protect\hust@item@square,
    leftmargin=2.0em, labelsep=0.6em
}
\setlist[itemize,2]{
    label=\protect\hust@item@circle,
    leftmargin=1.8em, labelsep=0.6em
}
\setlist[itemize,3]{
    label=\protect\hust@item@dash,
    leftmargin=1.8em, labelsep=0.6em
}
\end{verbatim}

\subsection{自定义有序列表样式（第 1045-1070 行）}

\begin{verbatim}
% 括号式：(1), (2), (3)
\newlist{enumpar}{enumerate}{1}
\setlist[enumpar]{label=(\arabic*)}

% 半括号式：1), 2), 3)
\newlist{enumbra}{enumerate}{1}
\setlist[enumbra]{label=\arabic*)}

% 圆圈数字
\newlist{enumcircled}{enumerate}{1}
\setlist[enumcircled]{
    label=\protect\hust@circled{\arabic*}
}
\end{verbatim}

% =========================================================
% 第十二章: 模块 18-19 - 实用功能与批注系统
% =========================================================
\newpage
\section{模块 18-19：实用功能与批注系统}

\subsection{目录树显示（第 1077-1078 行）}

\begin{verbatim}
\RequirePackage{dirtree}
\renewcommand{\DTstyle}{\ttfamily\small}
\end{verbatim}

\subsection{Todo/Fixme/Note 命令（第 1100-1139 行）}

可跨页的提示框：

\begin{verbatim}
\newcommand{\todo}[1]{%
    \begin{tcolorbox}[
        hust@todostyle,
        colback=orange!5,
        colframe=orange!80!black,
        title={TODO (待办)}
    ]
    #1
    \end{tcolorbox}%
}
\end{verbatim}

\subsection{PDF 元数据注入（第 1142-1149 行）}

\begin{verbatim}
\AtBeginDocument{
    \hypersetup{
        pdftitle={\@title},
        pdfauthor={\@author},
        pdfsubject={黑龙江科技大学实习报告/毕业设计},
        pdfcreator={HUST Report LaTeX Class}
    }
}
\end{verbatim}

\subsection{批注系统（第 1166-1284 行）}

使用 \texttt{pgfkeys} 解析选项：

\begin{verbatim}
\pgfkeys{
    /hustreview/.is family, /hustreview,
    style/.store in  = \hust@review@style,
    color/.store in  = \hust@review@color,
    side/.store in   = \hust@review@side,
    offset/.store in = \hust@review@offset
}

\NewDocumentCommand{\review}{ O{} m +m }{%
    \pgfkeys{/hustreview, #1}%
    % 根据样式渲染标记（下划线/波浪线等）
    % 使用 TikZ overlay 绘制侧边批注
}
\end{verbatim}

% =========================================================
% 第十四章: 实战效果展示（使用教程环境）
% =========================================================
\newpage
\section{实战效果展示}

本章使用模板自带的教程环境（\texttt{democode}、\texttt{demovert}、\texttt{codeonly}）来展示各功能的\textbf{实际渲染效果}，让读者直观理解每个命令和环境的用法。

\subsection{基础排版效果}

\subsubsection{字体命令演示}

\begin{democode}{三种预定义字体}
{\songti 这是宋体文本} \\
{\heiti 这是黑体文本} \\
{\kaiti 这是楷体文本}
\end{democode}

\subsubsection{字号命令演示}

\begin{democode}{常用字号对比}
{\zihao{3} 三号字} \\
{\zihao{4} 四号字} \\
{\zihao{-4} 小四号字} \\
{\zihao{5} 五号字}
\end{democode}

\subsubsection{文本修饰效果}

\begin{democode}{加粗与斜体}
普通文本、\textbf{加粗文本}、\textit{斜体文本}、\textbf{\textit{粗斜体}}
\end{democode}

% ---------------------------------------------------------
\subsection{列表环境效果}

\subsubsection{无序列表（自定义标记）}

\begin{demovert}{三级无序列表}
\begin{itemize}
    \item 一级列表项（实心方块）
    \begin{itemize}
        \item 二级列表项（灰色圆点）
        \begin{itemize}
            \item 三级列表项（短横线）
        \end{itemize}
    \end{itemize}
\end{itemize}
\end{demovert}

\subsubsection{有序列表}

\begin{demovert}{标准有序列表}
\begin{enumerate}
    \item 第一项内容
    \item 第二项内容
    \begin{enumerate}
        \item 子项目 (1)
        \item 子项目 (2)
    \end{enumerate}
    \item 第三项内容
\end{enumerate}
\end{demovert}

% ---------------------------------------------------------
\subsection{数学公式效果}

\subsubsection{行内公式}

\begin{democode}{行内数学公式}
质能方程 $E = mc^2$ 是物理学最著名的公式之一。
\end{democode}

\subsubsection{行间公式}

\begin{demovert}{带编号的行间公式}
牛顿第二定律表述为：
\begin{equation}
    F = ma
\end{equation}
其中 $F$ 为力，$m$ 为质量，$a$ 为加速度。
\end{demovert}

\subsubsection{物理符号（physics 宏包）}

\begin{demovert}{physics 宏包便捷命令}
偏导数：$\pdv{u}{t} + \pdv{u}{x} = 0$

矢量符号：$\vb{F} = m\vb{a}$

狄拉克符号：$\bra{\psi}\ket{\phi}$
\end{demovert}

% ---------------------------------------------------------
\subsection{提示框效果}

\subsubsection{注意框（蓝色主题）}

\begin{notice}{编译提醒}
本模板必须使用 XeLaTeX 编译，否则中文将无法正常显示！
\end{notice}

\subsubsection{定义框（绿色主题）}

\begin{definition}{边际消费倾向 (MPC)}
边际消费倾向是指每增加一单位可支配收入中用于消费的比例，计算公式为：
\[ MPC = \frac{\Delta C}{\Delta Y} \]
\end{definition}

\subsubsection{结论框（红色主题）}

\begin{conclusion}{核心发现}
通过本次调查研究，我们发现大学生消费行为主要受以下三个因素影响：
\begin{enumerate}
    \item 家庭经济状况
    \item 个人理财意识
    \item 社交消费需求
\end{enumerate}
\end{conclusion}

% ---------------------------------------------------------
\subsection{待办提示效果}

\todo{这是一个待办事项提醒框（橙色），用于标记需要完成的任务。}

\fixme{这是一个需要修复的问题标记（红色），用于标记错误或需要修正的内容。}

\note{这是一条备注信息（蓝色），用于补充说明或提供额外信息。}

% ---------------------------------------------------------
\subsection{API 参考卡片效果}

\apibox{makecover}{无参数}{生成标准报告封面页，自动读取 \cmd{title}、\cmd{author} 等预设信息。}

\apibox{maketaskbook}{无参数}{生成任务书页面，支持内容自动跨页。任务书内容通过 \cmd{renewcommand} 设置。}

\apibox{cnabstract}{\{关键词\}}{中文摘要环境，自动生成标题和关键词格式，并添加到目录。}

% ---------------------------------------------------------
\subsection{选项列表效果}

以下展示 \texttt{optionlist} 环境的实际效果：

\begin{optionlist}
    \option{win}{指定目标平台为 Windows，使用系统字体}
    \option{mac}{指定目标平台为 macOS，使用系统字体}
    \option{linux}{指定目标平台为 Linux，使用 Fandol 字体（默认）}
    \option{thesis}{启用毕业论文模式，应用特殊页眉设计}
\end{optionlist}

% ---------------------------------------------------------
\subsection{行内代码辅助命令}

\begin{democode}{cmd/env/opt 命令}
命令引用：\cmd{section} \\
环境引用：\env{itemize} \\
选项引用：\opt{win,thesis}
\end{democode}

% ---------------------------------------------------------
\subsection{表格效果}

\subsubsection{三线表}

\begin{demovert}{标准三线表}
\begin{table}[H]
\centering
\caption{消费类型统计}
\begin{tabular}{lcc}
\toprule
消费类型 & 占比 & 月均金额 \\
\midrule
餐饮 & 45\% & 900元 \\
学习用品 & 25\% & 500元 \\
娱乐 & 15\% & 300元 \\
其他 & 15\% & 300元 \\
\bottomrule
\end{tabular}
\end{table}
\end{demovert}

% ---------------------------------------------------------
\subsection{代码环境效果}

\subsubsection{codeonly 环境}

\begin{codeonly}{Python 数据处理示例}
import pandas as pd

def analyze_data(file_path):
    df = pd.read_csv(file_path)
    summary = df.describe()
    return summary
\end{codeonly}

% ---------------------------------------------------------
\subsection{图文混排效果}

\texttt{textfigure} 环境用于创建图文并排的布局。由于演示环境的嵌套限制，这里使用 \texttt{codeonly} 展示语法：

\begin{codeonly}{textfigure 环境语法}
% 语法：\begin{textfigure}[位置]{图片路径}{标题}[宽度比例]
% 位置参数：left = 左图右文，right = 右图左文（默认）
% 宽度比例：图片占行宽比例，默认 0.4

% 示例1：右图左文（默认）
\begin{textfigure}{figure/chart.png}{消费结构饼图}
    这是与图片并排显示的文字内容。
    文字会自动排列在图片一侧。
\end{textfigure}

% 示例2：左图右文 + 自定义宽度
\begin{textfigure}[left]{figure/trend.png}{趋势图}[0.5]
    从时间维度分析数据变化趋势，
    可以看到明显的增长模式。
\end{textfigure}
\end{codeonly}

\subsubsection{实际效果演示}

下面展示 \texttt{textfigure} 环境的实际渲染效果：

\begin{textfigure}{figure/example_chart.png}{消费结构分析}
    根据调查数据统计，大学生月均消费主要集中在以下几个方面：餐饮支出占比最高，约为 45\%；其次是学习用品支出，占 25\%；娱乐消费占 15\%；其他支出（如交通、通讯等）占 15\%。
    
    从数据可以看出，大学生的消费结构相对理性，基本生活需求占据主导地位。
\end{textfigure}

下面展示左图右文模式（\texttt{[left]} 参数）：

\begin{textfigure}[left]{figure/example_bar.png}{月度消费趋势}[0.45]
    从月度消费趋势来看，9月份因开学采购需求较高，消费水平达到峰值；10-11月趋于平稳；12月因期末考试和假期临近，娱乐消费有所下降。
    
    这一趋势与学期节奏高度吻合，反映了大学生消费行为的周期性特征。
\end{textfigure}

\textbf{参数说明}：
\begin{optionlist}
    \option{[位置]}{可选参数，\texttt{left}=左图右文，\texttt{right}=右图左文（默认）}
    \option{\{图片路径\}}{必填参数，图片文件的相对路径}
    \option{\{标题\}}{必填参数，图片的标题文字}
    \option{[宽度比例]}{可选参数，图片占行宽比例，范围 0-1，默认 0.4}
\end{optionlist}


% ---------------------------------------------------------
\subsection{并排多图效果}

\texttt{parallelfigures} 环境用于并排展示多张图片：

\begin{codeonly}{parallelfigures 环境语法}
% 语法：\begin{parallelfigures}{总标题}
%           \addfig[宽度]{路径}{子标题}
%           \addfig[宽度]{路径}{子标题}
%       \end{parallelfigures}

% 示例：三图并排
\begin{parallelfigures}{不同年级消费对比}
    \addfig[0.32]{figure/grade1.png}{大一}
    \addfig[0.32]{figure/grade2.png}{大二}
    \addfig[0.32]{figure/grade3.png}{大三}
\end{parallelfigures}

% 示例：两图并排（默认宽度 0.48）
\begin{parallelfigures}{性别消费差异}
    \addfig{figure/male.png}{男生}
    \addfig{figure/female.png}{女生}
\end{parallelfigures}
\end{codeonly}

\textbf{参数说明}：
\begin{itemize}
    \item \textbf{总标题}（必填）：整组图片的总标题
    \item \textbf{\textbackslash addfig} 命令：
    \begin{optionlist}
        \option{[宽度]}{可选，单图占行宽比例，默认 0.48}
        \option{\{路径\}}{必填，图片文件路径}
        \option{\{子标题\}}{必填，子图标题}
    \end{optionlist}
\end{itemize}

\subsubsection{实际效果演示}

下面使用 \texttt{parallelfigures} 环境展示两张示例图片的并排效果。注意子图编号自动采用 (a)、(b) 格式：

\begin{parallelfigures}{消费数据可视化示例}
    \addfig{figure/example_chart.png}{消费结构饼图}
    \addfig{figure/example_bar.png}{月度消费对比}
\end{parallelfigures}

如上图所示，主图编号采用"章节-序号"格式（如：图 14-1），子图编号采用 (a)、(b)、(c) 格式。引用时可以使用 \verb|\ref{fig:label}| 来引用主图或子图。

% ---------------------------------------------------------
\subsection{新增多图环境}

模板 v3.0 新增了两个多图环境，提供更灵活的排版方式：

\subsubsection{figurerow 环境 - 自动等宽行布局}

\begin{codeonly}{figurerow 环境语法}
% 语法：\begin{figurerow}{总标题}[列数]
%           \figitem{路径}{子标题}
%           \figitem{路径}{子标题}
%       \end{figurerow}

% 示例：3列布局
\begin{figurerow}{三张图并排}[3]
    \figitem{figure/img1.png}{子图A}
    \figitem{figure/img2.png}{子图B}
    \figitem{figure/img3.png}{子图C}
\end{figurerow}
\end{codeonly}

\subsubsection{figuregrid 环境 - 网格布局}

\begin{codeonly}{figuregrid 环境语法}
% 语法：支持2x2、2x3等任意网格布局
\begin{figuregrid}{2x2 网格图片}
    \gridrow{
        \gridfig{figure/a.png}{(a) 第一张}
        \gridfig{figure/b.png}{(b) 第二张}
    }
    \gridrow{
        \gridfig{figure/c.png}{(c) 第三张}
        \gridfig{figure/d.png}{(d) 第四张}
    }
\end{figuregrid}
\end{codeonly}

\textbf{子图编号说明}：
\begin{itemize}
    \item 主图编号：自动采用"章节-序号"格式，如 图 14-1
    \item 子图编号：自动采用 (a)、(b)、(c) 格式
    \item 引用语法：\verb|\ref{fig:main}| 引用主图，\verb|\ref{fig:sub}| 引用子图
\end{itemize}




模板自动设置"章节-序号"的编号格式，效果展示：

\begin{democode}{图表编号引用}
如\ref{fig:demo}所示... \\
根据\ref{tab:demo}可知... \\
由\ref{eq:demo}计算得...
\end{democode}

\textbf{实际编号格式}：图 14-1、表 14-1、式 (14-1)

% ---------------------------------------------------------
\subsection{交叉引用效果}

\begin{demovert}{自动前缀的引用}
使用 labelformat 设置后，引用自动添加前缀：
\begin{itemize}
    \item 图引用格式：图 X-Y
    \item 表引用格式：表 X-Y
    \item 公式引用格式：式 (X-Y)
\end{itemize}
\end{demovert}

% =========================================================
% 总结
% =========================================================
\newpage
\section*{总结}
\addcontentsline{toc}{section}{总结}

本文档对 \texttt{hustreport.cls} 类文件进行了全面分析，涵盖全部 20 个功能模块：

\begin{enumerate}
    \item 版本声明与类选项
    \item 按功能组织的宏包加载
    \item 用户信息接口
    \item 带回退链的智能字体配置
    \item 条件页面布局（报告模式 vs 论文模式）
    \item 章节标题与目录样式
    \item 浮动体优化与标题格式
    \item PDF 属性与超链接
    \item 特殊环境（摘要、参考文献、附录）
    \item TikZ 封面与任务书生成
    \item 高级图文混排布局
    \item 专业代码块与语法高亮
    \item 英文摘要环境
    \item 算法环境汉化
    \item 主题提示框（注意/定义/结论）
    \item 致谢环境
    \item TikZ 绘制的自定义列表标记
    \item 实用功能（目录树、待办、元数据）
    \item 高级批注与审阅系统
    \item 教程文档环境
\end{enumerate}

\textbf{第 14 章}通过 \texttt{democode}、\texttt{demovert}、\texttt{codeonly} 等教程环境，展示了模板各功能的\textbf{实际渲染效果}，帮助读者直观理解各命令和环境的用法。

\vspace{2cm}
\begin{center}
    \rule{0.4\textwidth}{0.5pt}\\[5pt]
    {\small 文档使用 hustreport.cls v3.0 生成}\\
    {\small \today}
\end{center}

\end{document}

